% ============================================================
\chapter{Vietoris-Rips, \v{C}ech, and Alpha Complexes}
\label{ch:complexes}
% ============================================================

Persistent homology needs a \emph{filtration} of simplicial complexes. But how does one build a complex from a point cloud? Three classical constructions dominate. The \emph{\v{C}ech complex} adds a simplex whenever the balls of radius $r$ centred at its vertices have a common intersection---it captures exactly the topology of the union of balls (nerve theorem). The \emph{Vietoris-Rips complex} is simpler: a simplex is present whenever all pairs of vertices are within distance $\leq 2r$---it is much easier to compute, but only approximates the \v{C}ech. The \emph{Alpha complex}, due to Edelsbrunner, restricts the \v{C}ech to Vorono\"i cells, producing a complex of linear size in the number of points.

\begin{intuition}
Constructing a simplicial complex from a point cloud is the fundamental
step of TDA. Three constructions dominate: Vietoris--Rips (simple but
large), \v{C}ech (exact but costly), and Alpha (small and exact in low
dimensions).
\end{intuition}

% ============================================================
\section{Vietoris--Rips Complex}
% ============================================================

\begin{definition}[Vietoris--Rips Complex]
Let $X = \{x_1, \ldots, x_n\}$ be a finite set in a metric space
$(M, d)$. The \emph{Vietoris--Rips complex} at parameter $r \geq 0$ is:
\[
  \mathrm{VR}(X, r) = \big\{
    \sigma \subseteq X : d(x_i, x_j) \leq r
    \;\text{for all}\; x_i, x_j \in \sigma
  \big\}.
\]
\end{definition}

\begin{remark}
The Vietoris--Rips complex is the \emph{clique complex} (or \emph{flag
complex}) of the proximity graph $G_r = (X, E_r)$ where
$\{x_i, x_j\} \in E_r$ iff $d(x_i, x_j) \leq r$. That is,
$\sigma \in \mathrm{VR}(X, r)$ iff all pairs of vertices of $\sigma$
are connected in $G_r$.
\end{remark}

\begin{proposition}[Properties]
\begin{enumerate}
  \item $\mathrm{VR}(X, r)$ is entirely determined by its $1$-skeleton
        (the graph $G_r$).
  \item $r \leq s \implies \mathrm{VR}(X, r) \subseteq \mathrm{VR}(X, s)$.
  \item The maximum dimension of $\mathrm{VR}(X, r)$ can reach $n - 1$.
\end{enumerate}
\end{proposition}

\begin{warning}[title=\textbf{Combinatorial Explosion}]
For $n$ points, the Vietoris--Rips complex can contain up to
$2^n - 1$ simplices. In practice, one truncates the maximum dimension
(typically $\leq 3$).
\end{warning}

% ============================================================
\section{\v{C}ech Complex}
% ============================================================

\begin{definition}[\v{C}ech Complex]
The \emph{\v{C}ech complex} at parameter $r$ is:
\[
  \check{C}(X, r) = \big\{
    \sigma \subseteq X : \bigcap_{x_i \in \sigma} B(x_i, r) \neq \emptyset
  \big\}
\]
where $B(x_i, r) = \{y \in M : d(x_i, y) \leq r\}$ is the closed ball
of center $x_i$ and radius $r$.
\end{definition}

\begin{theorem}[Nerve Theorem --- Borsuk, 1948]
Let $\mathcal{U} = \{U_i\}_{i \in I}$ be a finite cover of a
topological space $X$ by convex open sets. Then the \emph{nerve} of
$\mathcal{U}$:
\[
  \mathcal{N}(\mathcal{U}) = \big\{
    \sigma \subseteq I : \bigcap_{i \in \sigma} U_i \neq \emptyset
  \big\}
\]
has the same homotopy type as $\bigcup_{i \in I} U_i$.
\end{theorem}

\begin{corollary}
In $\R^d$, $\check{C}(X, r)$ has the same homotopy type as
$\bigcup_{i=1}^n B(x_i, r)$.
\end{corollary}

\begin{proposition}[Rips--\v{C}ech Relationship]
For any point cloud $X \subset \R^d$ and $r > 0$:
\[
  \check{C}(X, r) \subseteq \mathrm{VR}(X, 2r) \subseteq
  \check{C}\left(X, \frac{2r}{\sqrt{2}} \cdot \sqrt{\frac{d}{d+1}}\right).
\]
In particular: $\check{C}(X, r) \subseteq \mathrm{VR}(X, 2r)$.
\end{proposition}

% ============================================================
\section{Alpha Complex}
% ============================================================

\begin{definition}[Vorono\"{\i} Diagram]
The \emph{Vorono\"{\i} diagram} of $X \subset \R^d$ is the partition
of $\R^d$ into cells:
\[
  V_i = \{y \in \R^d : \norm{y - x_i} \leq \norm{y - x_j}
  \;\text{for all}\; j\}.
\]
\end{definition}

\begin{definition}[Delaunay Triangulation]
The \emph{Delaunay triangulation} $\mathrm{Del}(X)$ is the simplicial
complex dual to the Vorono\"{\i} diagram:
$\sigma = \{x_{i_0}, \ldots, x_{i_p}\} \in \mathrm{Del}(X)$ iff
$\bigcap_{k=0}^p V_{i_k} \neq \emptyset$.
\end{definition}

\begin{definition}[Alpha Complex]
The \emph{Alpha complex} at parameter $r$ is:
\[
  \mathrm{Alpha}(X, r) = \big\{
    \sigma \in \mathrm{Del}(X) :
    \bigcap_{x_i \in \sigma} (B(x_i, r) \cap V_i) \neq \emptyset
  \big\}.
\]
It is the subcomplex of the Delaunay triangulation formed by simplices
whose balls restricted to Vorono\"{\i} cells intersect.
\end{definition}

\begin{theorem}[Edelsbrunner, 1995]
The Alpha complex satisfies:
\begin{enumerate}
  \item $\mathrm{Alpha}(X, r)$ has the same homotopy type as
        $\bigcup_{i=1}^n B(x_i, r)$ (and hence as $\check{C}(X, r)$).
  \item $\dim(\mathrm{Alpha}(X, r)) \leq d$ (the ambient dimension).
  \item The total number of simplices is $O(n^{\lceil d/2 \rceil})$.
\end{enumerate}
\end{theorem}

\begin{keyformulas}[title=\textbf{Size Comparison}]
For $n$ points in $\R^d$:
\begin{center}
\begin{tabular}{lcc}
\toprule
\textbf{Complex} & \textbf{Max dimension} & \textbf{Nb simplices} \\
\midrule
Vietoris--Rips & $n - 1$ & $O(2^n)$ \\
\v{C}ech & $n - 1$ & $O(2^n)$ \\
Alpha & $d$ & $O(n^{\lceil d/2 \rceil})$ \\
\bottomrule
\end{tabular}
\end{center}
\end{keyformulas}

% ============================================================
\section{Practical Construction}
% ============================================================

\begin{codeblock}[title=\textbf{Rips Complex with GUDHI}]
\begin{minted}{python}
import gudhi
import numpy as np

# Point cloud
np.random.seed(42)
X = np.random.randn(50, 3)

# Rips complex
rips = gudhi.RipsComplex(points=X, max_edge_length=2.0)
st = rips.create_simplex_tree(max_dimension=3)

print(f"Number of simplices: {st.num_simplices()}")
print(f"Number of vertices: {st.num_vertices()}")
print(f"Dimension: {st.dimension()}")

# Persistence
pers = st.persistence()
gudhi.plot_persistence_diagram(pers)
\end{minted}
\end{codeblock}

\begin{codeblock}[title=\textbf{Alpha Complex with GUDHI}]
\begin{minted}{python}
import gudhi
import numpy as np

# 2D point cloud
np.random.seed(42)
theta = np.linspace(0, 2*np.pi, 80, endpoint=False)
X = np.column_stack([np.cos(theta), np.sin(theta)])
X += 0.1 * np.random.randn(*X.shape)

# Alpha complex
alpha = gudhi.AlphaComplex(points=X)
st = alpha.create_simplex_tree()

print(f"Number of simplices: {st.num_simplices()}")
print(f"Dimension: {st.dimension()}")

# Persistence
pers = st.persistence()
gudhi.plot_persistence_diagram(pers)
\end{minted}
\end{codeblock}

\begin{codeblock}[title=\textbf{\v{C}ech Complex Approximation}]
\begin{minted}{python}
import gudhi
import numpy as np

# 2D point cloud
X = np.random.randn(30, 2)

# The Alpha complex is homotopy-equivalent to Čech
# and much smaller
alpha = gudhi.AlphaComplex(points=X)
st = alpha.create_simplex_tree()

# Filtration values are squared radii of the
# smallest enclosing ball of each simplex
for simplex, filt in st.get_filtration():
    if len(simplex) <= 3:
        print(f"  {simplex}: radius^2 = {filt:.4f}")
    if filt > 0.5:
        break
\end{minted}
\end{codeblock}

% ============================================================
\section{Sparse Rips and Weighted Rips}
% ============================================================

\begin{definition}[Sparse Rips --- Sheehy, 2013]
The \emph{Sparse Rips complex} is an approximation of the Vietoris--Rips
complex containing $O(n)$ simplices per dimension, instead of $O(2^n)$.
For a parameter $\varepsilon > 0$, it produces a
$(1+\varepsilon)$-approximation of the persistence diagrams.
\end{definition}

\begin{codeblock}[title=\textbf{Sparse Rips with GUDHI}]
\begin{minted}{python}
import gudhi
import numpy as np

X = np.random.randn(500, 3)

# Sparse Rips (epsilon = 0.3)
sparse_rips = gudhi.RipsComplex(
    points=X, max_edge_length=2.0, sparse=0.3)
st = sparse_rips.create_simplex_tree(max_dimension=2)

print(f"Sparse Rips: {st.num_simplices()} simplices")

# Comparison with standard Rips
rips = gudhi.RipsComplex(points=X[:100], max_edge_length=2.0)
st_full = rips.create_simplex_tree(max_dimension=2)
print(f"Full Rips (100 pts): {st_full.num_simplices()} simplices")
\end{minted}
\end{codeblock}

% ============================================================
\section{Cubical Complexes}
% ============================================================

\begin{definition}[Cubical Complex]
A \emph{cubical complex} is the analogue of a simplicial complex where
the building blocks are cubes (products of intervals) instead of
simplices. It is particularly suited for grid data (images, 3D volumes).
\end{definition}

\begin{codeblock}[title=\textbf{Persistent Homology of an Image}]
\begin{minted}{python}
import gudhi
import numpy as np

# Create a simple image (ring)
x = np.linspace(-2, 2, 50)
y = np.linspace(-2, 2, 50)
XX, YY = np.meshgrid(x, y)
image = np.sqrt(XX**2 + YY**2)  # distance to origin

# Cubical complex
cc = gudhi.CubicalComplex(
    top_dimensional_cells=image.flatten(),
    dimensions=image.shape
)

# Persistence
pers = cc.persistence()
print("Image persistence:")
for dim, (b, d) in pers:
    if d - b > 0.5:
        print(f"  H_{dim}: ({b:.2f}, {d:.2f}), "
              f"pers={d-b:.2f}")
\end{minted}
\end{codeblock}

% ============================================================
\section{Choosing a Complex in Practice}
% ============================================================

\begin{algorithme}[title=\textbf{Complex Selection Guide}]
\begin{itemize}
  \item \textbf{Low-dimensional data} ($d \leq 3$): use the
        \textbf{Alpha} complex. It is small, exact, and fast to build.
  \item \textbf{High-dimensional or metric data}: use
        \textbf{Vietoris--Rips} (with \texttt{ripser} for speed).
        Truncate dimension to $2$ or $3$.
  \item \textbf{Grid data} (images, volumes): use a \textbf{cubical}
        complex.
  \item \textbf{Large number of points}: use \textbf{Sparse Rips} or
        subsample.
\end{itemize}
\end{algorithme}

% ============================================================
\section{Exercises}
% ============================================================

\begin{exercise}
For $4$ points at the vertices of a unit square, construct by hand the
complexes $\mathrm{VR}(X, r)$ and $\check{C}(X, r)$ for
$r = 0.5, 1.0, 1.5$. Verify the inclusion
$\check{C}(X, r) \subseteq \mathrm{VR}(X, 2r)$.
\end{exercise}

\begin{exercise}
Construct the Vorono\"{\i} diagram and Delaunay triangulation for $5$
random points in $\R^2$ (possibly with
\texttt{scipy.spatial.Delaunay}).
\end{exercise}

\begin{exercise}
Compare the computation times of Alpha and Rips complexes for point
clouds of $100$ to $1000$ points in dimensions $2$, $3$, and $10$.
\end{exercise}

\begin{exercise}
Show that the Vietoris--Rips complex of $n+1$ points in general
position in $\R^n$, for sufficiently large $r$, is the full simplex
$\Delta^n$.
\end{exercise}

\begin{exercise}
Compute the persistent homology of a binary image of your choice using
a cubical complex with GUDHI.
\end{exercise}
